\section{Background and Motivation}
\label{sec:bg}

\subsection{Expert Bytes on One Accelerator}
\label{sec:serving}

An MoE layer \(\ell\) holds \(E_{\ell}\) routed experts and sends each token to \(k_{\ell}\) of them.
Expert weights are the movable object.
Non-expert parameters, the KV cache, activations, and a conversion workspace are sized for a declared maximum batch and sequence length and are reserved outside \(B\).
The runtime's allocation problem is only the expert budget.
Host memory holds a packed low-precision copy and a packed high-precision copy of each expert.
The device holds at most one live copy.
A live copy occupies either \(m^{\mathrm{lo}}_{\ell}\) or \(m^{\mathrm{hi}}_{\ell}\) bytes, both measured from the packed representation rather than from a nominal bit width.
Promoting an expert therefore costs \(m^{\mathrm{hi}}_{\ell}-m^{\mathrm{lo}}_{\ell}\) extra device bytes if a low-precision copy was already resident, and it costs \(m^{\mathrm{hi}}_{\ell}\) bytes if the expert was absent.

Three published mechanisms cover different corners of this state space.
Uniform post-training quantization assigns every expert the same bit width offline~\cite{frantar2022gptq,lin2024awq,dettmers2022gptint8,xiao2023smoothquant}.
The configuration is feasible only when the chosen width's pool fits in \(B\), and it cannot move fidelity when the router's hot set changes.
Cross-layer prefetch issues transfers for the experts predicted over the next \(S\) layers on a side stream, overlapping them with the compute of layers \(\ell,\ldots,\ell{+}S{-}1\)~\cite{song2024promoe,eliseev2023fast,hwang2024pre}.
\(S\) does not change routing.
It changes only how many distinct expert bytes the cache must absorb.
Mixed-precision offload uses a cheaper representation on the miss path so that a cache fill moves fewer bytes~\cite{tang2024hobbit,huang2026dymoe}.
The bit width in that design is a property of the transfer, not a pin that is protected from eviction.

\subsection{Preliminary Measurements}
\label{sec:prelim}

The joint protocol in Section~\ref{sec:eval} repeats the following measurements on one runtime.
We report them here because they fix the quantities the model has to explain.
All three were taken on an NVIDIA RTX~A6000~\cite{nvidiaA6000datasheet}.
They are not comparable to each other as end-to-end system results: the density and hot-set traces use the model's native routing with experts forced to the low tier, and the horizon sweep uses a separate 20\,GB residency cap.

\paragraph{Activation density swings inside one request.}
We define the unique-expert ratio of a step as the number of distinct experts selected by top-\(k\) routing, divided by \(E_{\ell}\), averaged over routed layers.
The average is over five fixed prompt blocks.
Batch \(b\) uses the first \(b\) prompts of a 32-prompt block.
Prefill is the unique set over all non-padding tokens of a prompt capped at 2{,}048 tokens.
Decode is the unique set of the single next token, not an accumulation over a long generation.
Table~\ref{tab:density} shows the ratio on three models.
On Qwen3-30B at batch 32 the prefill ratio is 94.0\% and the next decode step is 26.0\%.
On Phi-3.5-MoE the prefill ratio is essentially the entire pool at every batch we measured.
A horizon chosen for the decode working set does not describe the prefill working set of the same request.

\begin{table}[t]
  \centering
  \caption{Preliminary unique-expert activation ratio (\%), RTX~A6000.
  Each cell averages routed layers and five prompt blocks.
  Prefill accumulates the prompt; decode is one subsequent token.
  The joint protocol repeats this table before any system comparison.}
  \label{tab:density}
  \small
  \begin{tabular}{@{}llrrrrrr@{}}
    \toprule
    Model & Stage & 1 & 2 & 4 & 8 & 16 & 32 \\
    \midrule
    Qwen3-30B & Decode & 6.3 & 7.9 & 11.8 & 16.1 & 20.8 & 26.0 \\
              & Prefill & 59.8 & 70.0 & 82.9 & 89.7 & 92.4 & 94.0 \\
    Qwen3-Next-80B & Decode & 1.9 & 3.7 & 6.7 & 14.9 & 25.4 & 39.2 \\
                   & Prefill & 24.6 & 35.3 & 48.8 & 67.1 & 78.9 & 86.2 \\
    Phi-3.5-MoE & Decode & 12.5 & 22.7 & 37.1 & 59.3 & 77.8 & 88.8 \\
                & Prefill & 98.2 & 98.8 & 99.4 & 99.8 & 99.9 & 100.0 \\
    \bottomrule
  \end{tabular}
\end{table}

\paragraph{The hot set is a property of the task.}
On Qwen3-30B, Layer~15, we counted every selected top-\(k\) dispatch under three full workloads (WikiText, GSM8K, and HumanEval) with the scheduler frozen and every expert held at low precision.
The ten most frequently selected experts were disjoint across the three workloads.
A high-precision assignment computed on one calibration corpus therefore pins the wrong experts on the other two tasks.
The slow loop has to be allowed to change identities.
It also has to be slow: a swap replaces a resident block and consumes migration bandwidth.

\paragraph{The horizon has a U-shaped latency curve when the pool does not fit.}
Under a 20\,GB device cap that forces expert absence, latency against \(S\) on three smaller MoE models was U-shaped.
Below the minimum, transfers do not finish inside the compute of the next \(S\) layers and the critical path accumulates exposed wait.
Above the minimum, the speculative set exceeds the cap, evictions increase, and later misses reload experts that a smaller horizon would have kept.
The location of the minimum differed by model on the same device.
Holding \(S\) fixed and increasing the batch from 1 to 32 increased the waiting component by up to \(2.5\times\).
At fixed batch size, requests with higher intra-batch embedding dispersion activated more distinct experts and showed up to \(300\%\) end-to-end latency variation against low-dispersion requests of the same length.
Effective host-to-device bandwidth across the devices used in that sweep spans roughly an order of magnitude, from a low-bandwidth discrete GPU to NVIDIA~H20 and Ascend~910B at about 128\,GB/s.
A horizon selected on one (device, batch, dispersion) triple is a different point on the curve for the others.

These three facts are the constraints on the model.
Density determines how many distinct experts a step demands.
Task identity determines which of those experts are worth a high-precision pin, on a slower timescale than a single layer.
Bandwidth, batch, and dispersion move the horizon that hides transfers, and they do so only when absence is unavoidable.
The rest of this section writes those constraints as a budget ratio and an overlap predicate.

\subsection{The Budget Ratio}
\label{sec:model}

Let \(P_{\mathrm{lo}}=\sum_{\ell} E_{\ell}\, m^{\mathrm{lo}}_{\ell}\) and \(P_{\mathrm{hi}}=\sum_{\ell} E_{\ell}\, m^{\mathrm{hi}}_{\ell}\).
Both sums are known once the packed representations have been materialized.
The deployment ratio is \(\PhiRatio=B/P_{\mathrm{lo}}\).
Staging for a publish-before-reclaim overlap is a constant \(M_{\mathrm{stage}}=K \max_{\ell,b} m^{b}_{\ell}\), with \(K\) the maximum number of in-flight transitions.
The controller spends \(B-M_{\mathrm{stage}}\).
We fold \(M_{\mathrm{stage}}\) into the definition of the usable budget from here on and continue to call it \(B\), so the inequalities below are stated against the usable budget.
Figure~\ref{fig:phase} is a schematic of the resulting regimes, not a measurement.

\begin{figure}[t]
  \centering
  \begin{tikzpicture}[
      font=\small,
      box/.style={draw, minimum height=1.15cm, inner sep=2pt, align=center},
      arr/.style={-{Stealth}, thick}
    ]
    \node[box, minimum width=3.3cm, fill=black!6] (p3) {Phase III\\ \(\PhiRatio<1\)\\ absent experts};
    \node[box, minimum width=3.5cm, fill=black!3, right=0pt of p3] (p2) {Phase II\\ \(1\le\PhiRatio<\rho\)\\ all resident at low};
    \node[box, minimum width=3.2cm, fill=white, right=0pt of p2] (p1) {Phase I\\ \(\PhiRatio\ge\rho\)\\ all resident at high};
    \node[below=0.15cm of p3, align=center] {tighter budget};
    \node[below=0.15cm of p1, align=center] {looser budget};
  \end{tikzpicture}
  \Description{Three adjacent boxes labeled Phase III, Phase II, and Phase I, ordered from a tighter budget to a looser budget.}
  \caption{Regimes of the expert-byte budget.
  \(\rho=P_{\mathrm{hi}}/P_{\mathrm{lo}}\).
  Phase~II spends leftover bytes on high-precision pins and sets the horizon to zero.
  Phase~III spends bytes first on the smallest feasible horizon, then on pins.
  The figure is the model's partition; Section~\ref{sec:eval} checks it against measured latency.}
  \label{fig:phase}
\end{figure}

\paragraph{Phase I.}
If \(B \ge P_{\mathrm{hi}}\), every expert fits at high precision.
The latency and quality costs of migration are both zero at the all-high assignment.
\systemname\ loads that assignment and disables both loops.

\paragraph{Phase II.}
If \(P_{\mathrm{lo}} \le B < P_{\mathrm{hi}}\), every expert fits at low precision and the all-high assignment does not fit.
Absence is unnecessary.
Any prefetch that moves an expert already resident adds a transfer and does not remove one.
The horizon that minimizes transfer latency is therefore \(S^{\star}=0\).
Replacing a low-precision resident with a high-precision resident costs \(m^{\mathrm{hi}}_{\ell}-m^{\mathrm{lo}}_{\ell}\) additional bytes.
When expert sizes are uniform within a layer and the same across layers, the maximum number of high-precision experts per layer satisfies
\begin{equation}
  n_{\mathrm{hi}}
  \;\le\;
  \left\lfloor
    \frac{B-P_{\mathrm{lo}}}{m^{\mathrm{hi}}-m^{\mathrm{lo}}}
    \cdot
    \frac{1}{L}
  \right\rfloor,
  \label{eq:phase2-nhi}
\end{equation}
with \(L\) the number of MoE layers, the floor applying after a proportional split.
The general case without that uniformity assumption is the integer knapsack that maximizes \(\sum_{\ell} n_{\mathrm{hi},\ell}\) subject to \(\sum_{\ell} n_{\mathrm{hi},\ell}(m^{\mathrm{hi}}_{\ell}-m^{\mathrm{lo}}_{\ell}) \le B-P_{\mathrm{lo}}\) and \(0 \le n_{\mathrm{hi},\ell} \le E_{\ell}\).
We use the proportional split in the implementation and record the knapsack as an optimality gap to measure, not as the online algorithm.
Identities inside the quota are the only online decision in this phase.
They are ranked by the slow hotness score of Section~\ref{sec:slow}.

As a numerical illustration of the algebra, not as a measured operating point: if \(m^{\mathrm{hi}}=4\,m^{\mathrm{lo}}\) and \(B=1.5\,P_{\mathrm{lo}}\), then \(\PhiRatio=1.5\) and the high-precision share of the pool is \((\PhiRatio-1)/(\rho-1)=1/6\).
About one sixth of the experts can be pinned high, and the rest stay resident at low precision.
A horizon controller that still evicts cold experts to ``make room'' destroys the floor that made \(S^{\star}=0\) feasible.

\paragraph{Phase III.}
If \(B < P_{\mathrm{lo}}\), a complete low-precision floor does not fit.
Some experts are absent, and a miss moves \(m\) bytes, where \(m\) is the size of the representation being fetched.
Let \(\hat{U}(S)\) be the set of experts the predictor expects to touch over the next \(S\) MoE layers, and let \(R\) be the set currently resident.
The miss set is \(\hat{U}(S)\setminus R\), and its transfer size is
\begin{equation}
  M_{\mathrm{miss}}(S)
  \;=\;
  \sum_{(\ell,e)\in \hat{U}(S)\setminus R} m_{\ell,e},
  \label{eq:miss-bytes}
\end{equation}
where \(m_{\ell,e}\) equals \(m^{\mathrm{hi}}_{\ell}\) if \((\ell,e)\) is a high-precision pin and \(m^{\mathrm{lo}}_{\ell}\) otherwise.
Pins are not misses.
A horizon \(S\) hides that transfer when two predicates hold simultaneously:
\begin{align}
  M_{\mathrm{miss}}(S) &\;\le\; S \cdot T_{\ell} \cdot C,
  \label{eq:overlap}\\
  \mathrm{bytes}(R \cup \hat{U}(S)) &\;\le\; B.
  \label{eq:fit}
\end{align}
Equation~\eqref{eq:overlap} is the overlap window.
Equation~\eqref{eq:fit} is cache pollution written as a hard constraint: a horizon whose speculative set does not fit is not a candidate, rather than a candidate with a penalty weight.
Define
\begin{equation}
  S^{\star}
  \;=\;
  \min\{\, S \in [S_{\min}, S_{\max}] : \eqref{eq:overlap}\ \mathrm{and}\ \eqref{eq:fit}\,\},
  \label{eq:sstar}
\end{equation}
and call the feasible set empty when no such \(S\) exists.
The minimum is the latency-optimal choice under a lexicographic rule: meet the overlap constraint, then leave the largest possible slack \(B-\mathrm{bytes}(R \cup \hat{U}(S^{\star}))\) for high-precision pins.
A larger feasible horizon would also hide the mean transfer and would consume bytes that could have raised \(n_{\mathrm{hi}}\).

The empty-set case is the dense sub-phase.
It is what Table~\ref{tab:density} suggests during prefill at high batch: \(\hat{U}(S)\) grows toward the full pool, \(M_{\mathrm{miss}}(S)\) exceeds both the overlap window and \(B\), and increasing \(S\) makes~\eqref{eq:fit} fail sooner.
The latency-minimizing action is then \(S^{\star}=S_{\min}\), with resident slots filled by the immediate layer's demand and then by the slow hotness order.
The non-empty case is the sparse sub-phase, typical of a small decode working set.
There \(S^{\star}\) is achievable, and every byte above \(\mathrm{bytes}(\hat{U}(S^{\star}))\) is slack.

Tail demand uses the same predicate.
Let \(M^{\mathrm{p95}}_{\mathrm{miss}}(S)\) be~\eqref{eq:miss-bytes} evaluated at the 95th percentile of the miss-set size observed over a sliding window of recent batches at this layer, with expert size \(m^{\mathrm{lo}}\) for unpinned experts.
The controller substitutes \(M^{\mathrm{p95}}_{\mathrm{miss}}\) for \(M_{\mathrm{miss}}\) in~\eqref{eq:overlap}.
The mean-demand horizon remains the quantity we plot against the model; the percentile horizon is what the runtime tracks.
Section~\ref{sec:eval} reports both, because a model that predicts the mean and misses the tail has not met a latency objective.

\subsection{What the Model Predicts}
\label{sec:predictions}

The regime rules are predictions about latency, not about task score.
Quality enters only as the use of slack.
We separate them so that a failure of the hotness proxy does not get mistaken for a failure of the overlap identity.

\begin{description}
  \item[P1.] In Phase~II the horizon that minimizes exposed expert-transfer latency is \(S=0\). A positive horizon increases migration bytes and does not reduce misses.
  \item[P2.] In Phase~III, when the feasible set is non-empty, \(S^{\star}\) is non-increasing in \(C\) and non-decreasing in the byte size \(m\) of a prefetched expert. Replacing a low-precision prefetch with a high-precision prefetch shifts \(S^{\star}\) up or empties the feasible set.
  \item[P3.] At fixed \(\PhiRatio<1\), a request can move from an empty feasible set during prefill to a non-empty set during decode. The minimizing action changes from ``fill the immediate resident set, hold \(S\) at \(S_{\min}\)'' to ``track \(S^{\star}\), then pin from the slack.''
  \item[P4.] Under the lexicographic rule, \(n_{\mathrm{hi}}\) is bounded by the slack after \(S^{\star}\) has been reserved. Raising \(S^{\star}\), raising \(m^{\mathrm{hi}}\), or entering the dense sub-phase all shrink that bound, including down to zero.
\end{description}

P1--P4 are the claims a joint re-measurement has to accept or reject.
They also state the boundary policies precisely.
A precision-only policy is the Phase~II optimum, including its \(n_{\mathrm{hi}}\) bound, applied at every \(\PhiRatio\).
When \(\PhiRatio<1\) that policy has no feasible full floor; an implementation must either exceed \(B\) or silently fall back to a tighter quantization.
A horizon-only policy is the Phase~III rule with \(m\) fixed and \(n_{\mathrm{hi}}\) fixed at zero, applied at every \(\PhiRatio\).
When \(\PhiRatio\ge 1\) it evicts experts the budget could have kept.
A naive composition runs both controllers and lets each assume it may use \(B\).
The composition violates~\eqref{eq:fit} whenever the pin quota and the speculative set are computed independently.
Section~\ref{sec:design} is the tracker that enforces one budget across both timescales.
Section~\ref{sec:eval} measures the four predictions before it measures that tracker against the boundary policies.

\subsection{Inputs the Model Does Not Observe}
\label{sec:model-limits}

The overlap identity uses a scalar \(C\).
Concurrent compute traffic reduces the bandwidth that a side-stream copy actually receives.
\systemname\ estimates \(C\) from completed transfers during overlap, not from an idle bandwidth benchmark, and Section~\ref{sec:eval} reports the error of \(S^{\star}\) under both estimators.
Dequantization and the launch of a second grouped GEMM, one per live precision in the layer, add device time that is not a host-to-device transfer.
Those times sit inside the measured \(T_{\ell}\).
They shrink the overlap window rather than appearing as a separate miss term.
The model does not predict task score.
Hotness is a proxy for which expert should receive a slack byte, and it can be wrong while P1--P3 still hold.
Kernel-launch overhead and proxy error are measured quantities in Section~\ref{sec:eval}, not terms hidden inside \(\PhiRatio\).
